%----------------------------------------------------------------------
%
%   The learned smoother for geometric multigrid on surfaces
%   Slide version of the 15th-meeting status report.
%
%   Same RUB/LKM beamer template as MS_Fanyang_15th: one folder per frame
%   under _Directory/, pulled in by \load, style files in _stylefiles/.
%   Sections mirror the six sections of the report one-for-one.
%
%----------------------------------------------------------------------

\pdfobjcompresslevel=1            			% solves error 131 in adobe reader on Windows
\documentclass[20pt]{beamer}
\usepackage{ifthen}							% needed for stylefiles

%-------------------------------------------------------------------
% set directory (in case slides are placed in different folder than main)
%-------------------------------------------------------------------
 \newcommand{\slides}{}

%-------------------------------------------------------------------
% lkm info input
%-------------------------------------------------------------------
\newcommand{\chairname}{Chair of Numerical Mathematics}	% chair name/ alternative: german
\newcommand{\profname}{Prof. Dr. Martin Kronbichler} 	% Prof. Name


%-------------------------------------------------------------------
% style files
%-------------------------------------------------------------------
% contrast color (alerted text, boxes, structure elements)
% true=green
% false=grey
\newboolean{contrastcolorstyle}
\setboolean{contrastcolorstyle}{true}
\usepackage{\slides_stylefiles/rublkm}
\usepackage{amsmath,amsfonts}
\usepackage{listings}
\usepackage{xcolor}
\usetikzlibrary{arrows.meta, positioning, shapes.geometric}
\usepackage{pgfplots}
\pgfplotsset{compat=1.18}
\usetikzlibrary{mindmap}
\usepackage{tikz}
% NOTE: no \bibliography here. The reference list is a literal thebibliography
% inside _Directory/Reference/slide.tex, so \cite resolves against it directly
% and bibtex is never run.

% definitions Daniel Balzani
% true  = italic bold variables in math mode
% false = roman bold variables in math mode
\newboolean{boldeqnitalic}						% boolean variable (defdb1):
\setboolean{boldeqnitalic}{true}
\usepackage{\slides_stylefiles/defdb1}

%-------------------------------------------------------------------
% notation of the report, so the slides and the report read alike
%-------------------------------------------------------------------
\newcommand{\Ac}{A_c}
\newcommand{\PiC}{\Pi_C}
\newcommand{\PiF}{\Pi_F}
\newcommand{\normA}[1]{\lVert #1\rVert_{A}}
\newcommand{\ETG}{E_{\mathrm{TG}}}
\newcommand{\Stheta}{S_{\theta}}
\newcommand{\Bnet}{B_{\theta}}
\newcommand{\rF}{\rho_F}
\newcommand{\rC}{\rho_C}
\newcommand{\muF}{\mu_F}
\newcommand{\muC}{\mu_C}

%-------------------------------------------------------------------
% colours and node styles of the two-level correction diagram
%-------------------------------------------------------------------
\definecolor{fine}{RGB}{35,91,160}
\definecolor{coarse}{RGB}{30,130,83}
\definecolor{corr}{RGB}{138,70,158}
\definecolor{smooth}{RGB}{184,105,31}
\tikzset{
 transfer/.style={-{Latex[length=2.7mm]},very thick},
 fn/.style={circle,fill=fine,inner sep=2.1pt},
 cn/.style={circle,fill=coarse,inner sep=2.6pt},
 cp/.style={circle,draw=corr,line width=1pt,minimum size=9pt,inner sep=0pt},
 meshlabel/.style={anchor=west,font=\small},
}

%-------------------------------------------------------------------
% load slides
%-------------------------------------------------------------------
\begin{document}
\titleslide{\load{preface}}

%%%%%%%%%%%%%%%%%%%%%%%%%%%%%%%%%%%%%%%%%%%%%%%%%%%%%%%%%%%%%%%%%%%%%
% Outline
%%%%%%%%%%%%%%%%%%%%%%%%%%%%%%%%%%%%%%%%%%%%%%%%%%%%%%%%%%%%%%%%%%%%%
\load{_Directory/Outline}

%%%%%%%%%%%%%%%%%%%%%%%%%%%%%%%%%%%%%%%%%%%%%%%%%%%%%%%%%%%%%%%%%%%%%
% 1. The question and the programme
%%%%%%%%%%%%%%%%%%%%%%%%%%%%%%%%%%%%%%%%%%%%%%%%%%%%%%%%%%%%%%%%%%%%%
\section{The Question and the Programme}
\load{_Directory/A1_Question}		% the working solver, the node, attribution, the three stages

%%%%%%%%%%%%%%%%%%%%%%%%%%%%%%%%%%%%%%%%%%%%%%%%%%%%%%%%%%%%%%%%%%%%%
% 2. The model problem
%%%%%%%%%%%%%%%%%%%%%%%%%%%%%%%%%%%%%%%%%%%%%%%%%%%%%%%%%%%%%%%%%%%%%
\section{The Model Problem}
\load{_Directory/B1_Model}			% the torus, the equation, the discretisation

%%%%%%%%%%%%%%%%%%%%%%%%%%%%%%%%%%%%%%%%%%%%%%%%%%%%%%%%%%%%%%%%%%%%%
% 3. Multigrid, and how the error propagates
%%%%%%%%%%%%%%%%%%%%%%%%%%%%%%%%%%%%%%%%%%%%%%%%%%%%%%%%%%%%%%%%%%%%%
\section{Multigrid and Error Propagation}
\load{_Directory/C1_TwoLevel}		% the hierarchy, the correction, C_h
\load{_Directory/C2_Projection}		% C_h = Pi_F exactly, E_TG
\load{_Directory/C3_Split}			% e = e_C + e_F, the direct sum, the extrapolation reading
\load{_Directory/C4_LearnedStep}	% where the network enters, one-sided vs symmetric, the leak

%%%%%%%%%%%%%%%%%%%%%%%%%%%%%%%%%%%%%%%%%%%%%%%%%%%%%%%%%%%%%%%%%%%%%
% 4. The learned smoother: the objective
%%%%%%%%%%%%%%%%%%%%%%%%%%%%%%%%%%%%%%%%%%%%%%%%%%%%%%%%%%%%%%%%%%%%%
\section{The Learned Smoother}
\load{_Directory/D1_Objectives}		% the three objectives
\load{_Directory/D2_Loss}			% which one, the constraints, what is left out

%%%%%%%%%%%%%%%%%%%%%%%%%%%%%%%%%%%%%%%%%%%%%%%%%%%%%%%%%%%%%%%%%%%%%
% 5. Numerical experiments
%%%%%%%%%%%%%%%%%%%%%%%%%%%%%%%%%%%%%%%%%%%%%%%%%%%%%%%%%%%%%%%%%%%%%
\section{Numerical Experiments}
\load{_Directory/E1_Classical}		% the baseline is not weak
\load{_Directory/E2_Setting}		% the Stage I experiment, and what is asserted
\load{_Directory/E3_Diagnostics}	% trunk floor, rank floor
\load{_Directory/E4_Measurement}	% the table and the answer
\load{_Directory/E5_Defect}			% per mechanism, and the defect

%%%%%%%%%%%%%%%%%%%%%%%%%%%%%%%%%%%%%%%%%%%%%%%%%%%%%%%%%%%%%%%%%%%%%
% 6. Conclusions
%%%%%%%%%%%%%%%%%%%%%%%%%%%%%%%%%%%%%%%%%%%%%%%%%%%%%%%%%%%%%%%%%%%%%
\section{Conclusions, Protocol, and Next Steps}
\load{_Directory/F1_Conclusion}

%%%%%%%%%%%%%%%%%%%%%%%%%%%%%%%%%%%%%%%%%%%%%%%%%%%%%%%%%%%%%%%%%%%%%
% References
%%%%%%%%%%%%%%%%%%%%%%%%%%%%%%%%%%%%%%%%%%%%%%%%%%%%%%%%%%%%%%%%%%%%%
\load{_Directory/Reference}

\end{document}
